% =====================================================================
%  CONJURA SOLUTION / PARTIAL PROOF
%  Extraction and decomposition for two split unpredictable sources.
%  Requires conjura-solution.cls and conjura-conjecture.cls here.
% =====================================================================
\documentclass{conjura-solution}

\runninghead{EXTRACTION FOR SPLIT UNPREDICTABLE SOURCES}

\newcommand{\Fun}{\mathsf{Fun}}
\newcommand{\Adv}{\mathsf{Adv}}
\newcommand{\Real}{\mathsf{Real}}
\newcommand{\Dec}{\mathsf{Dec}}
\newcommand{\Prob}[1]{\Pr\!\left[#1\right]}
\newcommand{\Exp}{\mathbb{E}}
\newcommand{\bits}{\{0,1\}}
\newcommand{\vx}{\mathbf{x}}
\newcommand{\vz}{\mathbf{z}}
\newcommand{\vu}{\mathbf{u}}
\newcommand{\ind}{\mathbb{1}}
\newcommand{\Unif}{\mathrm{Unif}}
\newcommand{\supp}{\mathrm{supp}}
\newcommand{\kna}{\kappa^{\mathrm{na}}}

\begin{document}

\cjkicker{CONJURA \textperiodcentered{} PARTIAL SOLUTION}
\cjtitle{Extraction for Split Unpredictable Sources, and Decomposition on a Region}
\cjresolves{\emph{Decomposition for Split Unpredictable Sources} (see
  \texttt{main.tex} in this folder). \textbf{Partially.} Three unconditional
  extraction bounds are proved outright, one of them free of the query budget
  altogether and tight in its dependence on $M$; the conjecture itself is proved,
  for \emph{every} observer, on a region of $P$, $q$ and $M$, and remains open in
  general.}

\vspace{0.4em}

\section*{Reader's guide}

\noindent\textbf{What is proved.} Write $\kna(q)$ for the extraction advantage
restricted to observers whose query positions, for each fixed coin string, do not
depend on the challenge value, and
$\mu(s):=\min\bigl(s\delta,\,2(s\delta^{2})^{1/3},\,1\bigr)$. Put
$\Lambda:=\min\bigl\{5\sqrt{\sigma'\delta},\ \sqrt{2\delta\ln(eN)}+4\delta\sqrt{\sigma'}\bigr\}$,
the second arm being the smaller once $\delta\le0.85$.

\begin{itemize}
\item[(A)] $\kna(q)\le\Lambda+\mu(q)$ for every $q$, and $\kappa(0)\le\Lambda$
  against \emph{all} observers. Unconditional (Theorem~\ref{thm:A}). This settles
  the statement's assertion that queries whose positions ignore the challenge
  value are almost free at every budget: their cost is the additive
  $\mu(q)\le q\delta$, with no multiplier on the leading term and no dependence
  on $M$.
\item[(B)] $\kappa(q)\le\Lambda+\mu(\min(qM,N^{2}))$ for every $q$ and every
  observer. Unconditional (Theorem~\ref{thm:B}).
\item[(D)] $\kappa(q)\le5\sqrt{\sigma'\delta}+2\delta\sqrt M$ for every $q$ and
  every observer --- indeed for observers of \emph{unbounded} query count.
  Unconditional (Theorem~\ref{thm:D}). Its display mentions neither $q$ nor
  $\sigma$, and the collected form above is free of $q$ and retains $\sigma$ only
  inside $\sigma'$; for every $q\ge1$ and $M\ge4$ it is the smaller of (B) and (D)
  wherever either says anything at all (Corollary~\ref{cor:D}).
\item[(C)] An explicit family, chosen before $q$, indexed by the oracle and the
  leakage, and consistent with the oracle, with
  $\Adv_{\mathcal Y,D}\le\varepsilon(D)+2(\sigma'+\log\gamma^{-1})q/P+\gamma+P\delta+q\delta$
  (Theorem~\ref{thm:C}).
\item[(E)] Hence the conjecture, with the absolute constants $c=2$, $C=8$, for
  every $P\le\sqrt{\sigma'/\delta}$ and every $q$ with $q^{+}\delta\le\sigma'$,
  and for \emph{every} observer once $M\le\sigma'q^{+}/(4\delta)$
  (Theorem~\ref{thm:cprime}). On that region no hypothesis whatever on how the
  observer's probes depend on the challenge value is needed.
\item[(F)] The term $2\delta\sqrt M$ of (D) is tight to within a factor $12$,
  already for observers whose query positions ignore the challenge
  (Proposition~\ref{prop:F}). No bound free of $q$ can do better in $M$.
\end{itemize}

\noindent\textbf{What is not proved.} The conjecture in full. Two obstructions
survive, stated as mechanisms in \S\ref{sec:barriers}: challenge-steered probes
at $M$ \emph{beyond} $\sigma'q^{+}/(4\delta)$, and large $P$ with growing $q$. The
statement's own remark already names the second as a separate problem.

\medskip

\noindent\textbf{Where to be most sceptical.} Two rounds of blind review found
two serious defects, both since repaired, and both of the same species --- a
claim left \emph{underdetermined} rather than false. They were at
Theorem~\ref{thm:cprime}'s first proof line (which quantity the cited corollary
actually bounds) and in Definition~\ref{def:res} (whether challenge resolution is
a condition on the law of the query positions or on each coin string). Those two
places are where a third error, if there is one, is most likely to be. Since then
the document has grown by Lemma~\ref{lem:mom}, Lemma~\ref{lem:disc2},
Theorem~\ref{thm:D}, branch~1 of Theorem~\ref{thm:cprime} and
Proposition~\ref{prop:F}, none of which has been reviewed at all. The review
status is recorded honestly in \S\ref{sec:review}.

\medskip

\noindent\textbf{Notation} is that of the statement document and is not repeated:
$\Fun=\{f:[N]\times[N]\to[M]\}$, $H\sample\Fun$, split $\delta$-unpredictable
sources $(S_1,S_2)$ with leakage bounded by $\sigma_1,\sigma_2$,
$\sigma=\sigma_1+\sigma_2$, $\sigma'=\sigma+2\log N$, $q^{+}=q+1$, and the
experiments $\Real,\Dec,\Real_0,\Dec_0$ with $\Adv_{\mathcal Y,D}$ and
$\kappa(q)$. Throughout $N\ge2$ and $M\ge2$; logarithms are base two. Write
$\mathcal Z:=\bits^{\le\sigma_1}\times\bits^{\le\sigma_2}$, so
$|\mathcal Z|=(2^{\sigma_1+1}-1)(2^{\sigma_2+1}-1)<2^{\sigma+2}$.

\section{The two structural facts}

Everything below rests on two observations, which together turn the conjecture
into a statement about the discrepancy of a random matrix on combinatorial
rectangles.

\begin{lemma}[derandomisation]\label{lem:derand}
Fix one of the two paired experiments, $(\Real,\Real_0)$ or $(\Real,\Dec)$. For
every $q$ and every $q$-query challenge-oblivious $D$ there is a deterministic
such $D'$ with \emph{(i)} the advantage \emph{in that pair} no larger for $D$ than
for $D'$ --- that is, $\bigl|\Prob{\Real=1}-\Prob{\Real_0=1}\bigr|$ in the first
case and $\Adv_{\mathcal Y,D}\le\Adv_{\mathcal Y,D'}$ in the second --- and
\emph{(ii)} $D'$ of the same challenge resolution as $D$
(Definition~\ref{def:res}). Consequently the suprema defining $\kappa(q)$ and
$\kna(q)$ are both attained on deterministic observers, the latter on
deterministic observers of resolution~$1$.
\end{lemma}

The maximiser is chosen per pair and the two need not coincide: the lemma does
\emph{not} supply one $D'$ dominating both experiments, and nothing here asks it
to. Every appeal to it --- in Theorems~\ref{thm:A}, \ref{thm:B} and~\ref{thm:D}
--- is to the pair $(\Real,\Real_0)$, and Theorem~\ref{thm:C} is stated for a
fixed $D$ with no supremum taken, so $\Adv_{\mathcal Y,D}$ is never bounded by
passing to a derandomised observer.

\begin{proof}
Let $\rho$ be $D$'s coins and $D_\rho$ the deterministic strategy it then runs.
All of $\Fun,[N],[M],\mathcal Z$ are finite and $D$ makes at most $q$ queries, so
there are finitely many deterministic $q$-query challenge-oblivious strategies.
Each quantity at issue is $|\Exp_\rho[\alpha(\rho)]|$ with
$\alpha(\rho)=\Prob{\mathsf E_1=1\mid\rho}-\Prob{\mathsf E_0=1\mid\rho}$, the
coins being independent of every other draw and the paired experiments differing
in no other component; so it is at most $\max_\rho|\alpha(\rho)|$, attained at
some $\rho^{*}$ depending on the pair fixed at the outset. Take
$D':=D_{\rho^{*}}$, which is $q$-query and challenge-oblivious. It inherits $D$'s challenge resolution: the definition
quantifies over every coin string, so if $D^{f}_{\rho}(v,\zeta)$ and
$D^{f}_{\rho}(v',\zeta)$ issue the same query positions for every $\rho$ whenever
$h(v)=h(v')$, they do so in particular at $\rho^{*}$, which is what $D'$ runs.
The argument is identical for whichever paired experiment was fixed, and is run
separately for each.
\end{proof}

Once $D$ is deterministic and its oracle is fixed to $f$, its output is a
function of its challenge input alone; write $\theta:[M]\to\bits$ for that
function and $p_\theta:=|\theta^{-1}(1)|/M$. That the test is \emph{named by
short data} is the crux of Lemma~\ref{lem:disc}.

\begin{lemma}[posterior product structure]\label{lem:prod}
For every $(f,\zeta)$ in the support of $(H,\vz)$ the conditional law of $\vx$ is
the product $\Pi_{f,\zeta}=\pi^{1}_{f,\zeta_1}\otimes\pi^{2}_{f,\zeta_2}$, where
$\pi^{i}_{f,\zeta_i}(u)=\Prob{x_i=u\mid H=f,\ z_i=\zeta_i}$. Writing
$m_i:=\|\pi^{i}_{f,\zeta_i}\|_\infty$:
\begin{enumerate}
\item $\max_{\vu\in[N]^2}\Pi_{f,\zeta}(\vu)=m_1m_2$;
\item $\Exp[m_i]\le\delta$ for $i=1,2$;
\item $\Exp[m_1m_2]\le\delta$;
\item $\Exp\bigl[\sqrt{m_1m_2}\bigr]\le\delta$;
\item $\Prob{m_1m_2>\vartheta}\le\delta/\sqrt\vartheta$ for every
  $\vartheta\in(0,1]$.
\end{enumerate}
\end{lemma}

\begin{proof}
The sources run on independent private coins and on the same $H$, so conditioned
on $H=f$ their output pairs are independent:
$\Prob{x_1=u_1,z_1=\zeta_1,x_2=u_2,z_2=\zeta_2\mid f}=\prod_i\Prob{x_i=u_i,z_i=\zeta_i\mid f}$.
Summing over $u_1,u_2$ gives
$\Prob{\vz=\zeta\mid f}=\prod_i\Prob{z_i=\zeta_i\mid f}>0$, and dividing gives the
product form. Item~1 is the fact that the maximum of a product measure is the
product of the maxima.

For item~2, let $\mathsf P$ be the predictor that, on input $\vz$ and with oracle
access to $H$, reads all of $H$, computes $\pi^{i}_{H,z_i}$ --- it may, the
sources being fixed and predictors being arbitrary, possibly randomised,
functions of their inputs, as the statement document's own proof of its
fixed-set lemma instructs --- and outputs a maximiser. Then
$\Prob{\mathsf P^{H}(\vz)=x_i}=\Exp_{(f,\zeta)}[m_i(f,\zeta)]$, which
unpredictability bounds by $\delta$. Item~3 follows from $m_2\le1$ and item~2.
Item~4 is Cauchy--Schwarz:
$\Exp[\sqrt{m_1}\cdot\sqrt{m_2}]\le\sqrt{\Exp[m_1]}\sqrt{\Exp[m_2]}\le\delta$, by
item~2 applied to each coordinate separately. Item~5 is Markov applied to the
non-negative variable $\sqrt{m_1m_2}$, together with item~4.
\end{proof}

Item~1 is where the second source is spent: a \emph{single cell} of $[N]^2$
carries mass $m_1m_2$, a product of two independently small numbers, whereas one
unpredictable source over $[N]^2$ gives only $\delta$ per cell.

\begin{remark}[{$\Exp[m_1m_2]\le\delta^{2}$ is false}]\label{rem:corr}
This matters, because it is what forces the exponent $1/3$ in
Lemma~\ref{lem:mass}. Let $2\le M\le N$, $\delta:=1/M$, let $E$ be the event
$f(1,1)=1$, of probability $\delta$, and let both sources output $x_i:=1$ if $E$
holds and $x_i$ uniform on $[N]$ otherwise, leaking the empty string. The coins
are independent, so the pair is split, and
$\Exp[m_i]=\delta+(1-\delta)/N\le2\delta$, so the pair is
$2\delta$-unpredictable; yet $m_1m_2=1$ on $E$, so $\Exp[m_1m_2]\ge\delta$.
Splitness constrains the coins, not the two sources' common dependence on $f$.
\end{remark}

\begin{remark}[the second source, spent twice]\label{rem:cs}
Item~4 of Lemma~\ref{lem:prod} does \emph{not} follow from item~3 by Jensen, which
runs the other way, and it is a factor $\sqrt\delta$ stronger than the
$\sqrt{\Exp[m_1m_2]}\le\sqrt\delta$ that concavity gives. It is where the second
source's independence is spent a second time, the first being item~1, and it is
what turns the union-bound budget of Lemma~\ref{lem:disc2} into $\delta\sqrt M$
rather than $\sqrt{M\delta}$ --- the difference between a useful bound and a
useless one. It is an equality on two structurally different families: on the one
in Remark~\ref{rem:corr}, where $m_1m_2$ has a heavy tail and
$\Exp[\sqrt{m_1m_2}]=\delta+(1-\delta)/N=\Exp[m_i]$ exactly, and on the flat family
of Proposition~\ref{prop:F}, where $m_1=m_2$ is deterministic.
\end{remark}

\section{Flattening and rectangle discrepancy}

\begin{lemma}[flattening]\label{lem:flat}
Let $\pi$ be a distribution on $[N]$ with $\|\pi\|_\infty=m$ and
$k:=\lfloor1/m\rfloor$. Then $1\le k\le N$, $1/k\le2m$, and $\pi$ is a convex
combination of distributions each uniform on a subset of size exactly $k$.
Consequently, for every $G:[N]^{2}\to[-1,1]$ and every $\pi^{1},\pi^{2}$ with
$\|\pi^{i}\|_\infty=m_i$, $k_i:=\lfloor1/m_i\rfloor$,
\[
  \bigl|\Exp_{\pi^{1}\otimes\pi^{2}}[G]\bigr|\ \le\
  \max\bigl\{\ \bigl|\Exp_{\Unif(B_1)\otimes\Unif(B_2)}[G]\bigr|\ :\ |B_i|=k_i\ \bigr\}.
\]
\end{lemma}

\begin{proof}
$m\ge1/N$ gives $k\le N$; $m\le1$ gives $k\ge1$. For $x\ge1$,
$\lfloor x\rfloor\ge x/2$ (for $1\le x<2$ because $\lfloor x\rfloor=1\ge x/2$;
for $x\ge2$ because $\lfloor x\rfloor>x-1\ge x/2$), so with $x=1/m$,
$1/k\le2m$. If $k=1$ the claimed decomposition is the identity
$\pi=\sum_u\pi(u)\delta_u$, a convex combination of distributions uniform on
singletons, and no external result is needed; assume $k\ge2$. Since
$\|\pi\|_\infty\le1/k=2^{-\log k}$, $\log k>0$ and $2^{\log k}=k\in\mathbb N$,
the flat-decomposition lemma [CFHS, Lemma~3] applies and writes $\pi$ as a convex
combination of distributions uniform on $k$-subsets. The displayed inequality
follows because $(\rho^1,\rho^2)\mapsto\Exp_{\rho^1\otimes\rho^2}[G]$ is
bilinear, so substituting both decompositions expresses the left side as a convex
combination of the quantities on the right.
\end{proof}

Note $G$ is arbitrary: it will carry an indicator $\ind[\vu\notin\mathcal S]$ that
does \emph{not} factor over the coordinates. That is harmless, because the linear
structure used lives in $(\pi^1,\pi^2)$, not in $G$.

Both concentration bounds below are one Hoeffding estimate with two different
union bounds, and both are stated through a single quantity. For $c>0$ put
$L:=\ln(eN)$ and
\[
  t_c(k_1,k_2):=\sqrt{\frac{k_1\ln\frac{eN}{k_1}+k_2\ln\frac{eN}{k_2}+c}{2k_1k_2}}\,,
\]
where $c$ is the logarithm of the number of tuples the union runs over.

\begin{lemma}[the expectation of $t_c$]\label{lem:mom}
For every $c>0$, pointwise at $k_i=\lfloor1/m_i\rfloor$,
$t_c(k_1,k_2)^{2}\le L(m_1+m_2)+2c\,m_1m_2$, and therefore
\[
  \Exp\bigl[t_c(k_1,k_2)\bigr]\ \le\
  \min\Bigl\{\ \sqrt{2\delta(L+c)}\ ,\ \ \sqrt{2L\delta}+\delta\sqrt{2c}\ \Bigr\}.
\]
\end{lemma}

\begin{proof}
Expanding,
$t_c^{2}=\frac{\ln(eN/k_1)}{2k_2}+\frac{\ln(eN/k_2)}{2k_1}+\frac{c}{2k_1k_2}$; each
$k_i\ge1$ gives $\ln(eN/k_i)\le L$, and $1/k_i\le2m_i$ from Lemma~\ref{lem:flat}
gives $\frac1{2k_i}\le m_i$ and $\frac1{2k_1k_2}\le2m_1m_2$, which is the pointwise
bound. For the first arm, concavity of $\sqrt\cdot$ applied to the whole expression
with Lemma~\ref{lem:prod}(2),(3) gives
$\Exp[t_c]\le\sqrt{L(\Exp m_1+\Exp m_2)+2c\Exp[m_1m_2]}\le\sqrt{2\delta(L+c)}$. For
the second, $\sqrt{a+b}\le\sqrt a+\sqrt b$ splits the pointwise bound first;
concavity and Lemma~\ref{lem:prod}(2) give
$\Exp[\sqrt{L(m_1+m_2)}]\le\sqrt{2L\delta}$, and Lemma~\ref{lem:prod}(4) gives
$\sqrt{2c}\,\Exp[\sqrt{m_1m_2}]\le\delta\sqrt{2c}$.
\end{proof}

Splitting before averaging is what the second arm buys: it pays $\delta\sqrt{2c}$
where the first pays $\sqrt{2c\delta}$. At $c=C_0$ below this is a sharpening of a
constant; at $c=C_1$, where $c$ carries the alphabet size, it decides whether the
bound says anything.

\begin{definition}[revealing rule]\label{def:rr}
A \emph{revealing rule of budget $s$} is a family
$(\mathcal A_\zeta)_{\zeta\in\mathcal Z}$ of procedures, where $\mathcal A_\zeta$
inspects the coordinates of $f$ one at a time --- the next a function of $\zeta$
and the values already seen --- halts having inspected a set
$\mathcal S(f,\zeta)$ with $|\mathcal S(f,\zeta)|\le s$, and outputs a test
$\theta_{\zeta,f}:[M]\to\bits$ that is a function of $\zeta$ and the inspected
values alone.
\end{definition}

\begin{lemma}[rectangle discrepancy against leakage-named tests]\label{lem:disc}
Let $\gamma_0\in(0,1)$ and let $(\mathcal A_\zeta)$ be a revealing rule. Put
$C_0:=(\sigma+2)\ln2+\ln\frac{4N^{2}}{\gamma_0}$. Then with probability at least
$1-\gamma_0$ over $f\sample\Fun$, simultaneously
for every $\zeta\in\mathcal Z$, all $k_1,k_2\in[N]$ and all $B_i\subseteq[N]$
with $|B_i|=k_i$, writing $R:=B_1\times B_2$,
\begin{equation}\label{eq:disc}
  \Bigl|\ \frac{1}{k_1k_2}\!\!\sum_{u\in R\setminus\mathcal S(f,\zeta)}\!\!
  \bigl(\theta_{\zeta,f}(f(u))-p_{\theta_{\zeta,f}}\bigr)\Bigr|\ \le\ t_{C_0}(k_1,k_2).
\end{equation}
\end{lemma}

\begin{proof}
\emph{Conditional uniformity off the revealed set.} Fix $\zeta$. For
$S_0\subseteq[N]^{2}$ and $w\in[M]^{S_0}$ the event
$\mathcal T_{S_0,w}:=\{\mathcal S(f,\zeta)=S_0\wedge f|_{S_0}=w\}$ is measurable
with respect to $f|_{S_0}$: replaying $\mathcal A_\zeta$ on $w$ reproduces its
whole inspection sequence. The coordinates of $f$ being i.i.d.\ uniform,
conditioning on such an event leaves $f|_{[N]^{2}\setminus S_0}$ i.i.d.\ uniform;
and $\theta_{\zeta,f}$, $p_{\theta_{\zeta,f}}$ are constant on
$\mathcal T_{S_0,w}$.

\emph{One tuple.} Fix also $k_1,k_2,B_1,B_2$ and condition on
$\mathcal T_{S_0,w}$. Write $\theta,p_\theta$ for the constants and
$n_0:=|R\setminus S_0|\le k_1k_2$. The variables $\theta(f(u))-p_\theta$,
$u\in R\setminus S_0$, are independent, lie in an interval of length $1$, and have
mean $0$, since $f(u)$ is uniform on $[M]$ and
$\Prob{\theta(f(u))=1}=p_\theta$ exactly. Hoeffding's inequality with
$\lambda:=k_1k_2\,t_{C_0}(k_1,k_2)$ gives conditional failure probability at most
$2\exp(-2\lambda^{2}/n_0)\le2\exp(-2t_{C_0}^{2}k_1k_2)$, uniformly in $(S_0,w)$,
hence also unconditionally.

\emph{Union bound.} By the definition of $t_c$,
$2\exp(-2t_{C_0}^{2}k_1k_2)=2\,e^{-k_1\ln\frac{eN}{k_1}}e^{-k_2\ln\frac{eN}{k_2}}2^{-(\sigma+2)}\frac{\gamma_0}{4N^{2}}$.
There are $|\mathcal Z|\binom{N}{k_1}\binom{N}{k_2}$ tuples with the given
$(k_1,k_2)$, and $\binom{N}{k}\le(eN/k)^{k}$ [CFHS, Lemma~4.3] while
$|\mathcal Z|<2^{\sigma+2}$. Multiplying, all three exponential factors cancel
and the contribution of the pair $(k_1,k_2)$ is at most
$2\gamma_0/(4N^{2})=\gamma_0/(2N^{2})$; summing over the $N^{2}$ pairs gives
$\gamma_0/2\le\gamma_0$.
\end{proof}

\begin{lemma}[rectangle discrepancy against \emph{all} tests]\label{lem:disc2}
Let $\gamma_0\in(0,1)$ and put $C_1:=M\ln2+\ln\frac{4N^{2}}{\gamma_0}$. Then with
probability at least $1-\gamma_0$ over $f\sample\Fun$, simultaneously for all
$k_1,k_2\in[N]$, all $B_i\subseteq[N]$ with $|B_i|=k_i$, and \emph{every} function
$\theta:[M]\to\bits$, writing $R:=B_1\times B_2$,
\[
  \Bigl|\ \frac{1}{k_1k_2}\sum_{\vu\in R}\bigl(\theta(f(\vu))-p_\theta\bigr)\Bigr|
  \ \le\ t_{C_1}(k_1,k_2).
\]
Equivalently, every rectangle's pushforward $\nu_R:=f_{*}\Unif(R)$ satisfies
$\mathrm{SD}\bigl(\nu_R,\Unif([M])\bigr)\le t_{C_1}(|B_1|,|B_2|)$.
\end{lemma}

\begin{proof}
Fix $k_1,k_2$, sets $B_i$ with $|B_i|=k_i$, and a function $\theta$ --- all fixed
before $f$ is drawn. The variables $\theta(f(\vu))-p_\theta$, $\vu\in R$, are
independent, the coordinates of $f$ being i.i.d.\ uniform; each has mean exactly
$0$, since $\Prob{\theta(f(\vu))=1}=|\theta^{-1}(1)|/M=p_\theta$; and each lies in
$\{-p_\theta,1-p_\theta\}$, an interval of length $1$. Hoeffding with
$\lambda:=k_1k_2\,t_{C_1}$ gives failure probability at most
$2\exp(-2t_{C_1}^{2}k_1k_2)$, which by the definition of $t_c$ equals
$2(eN/k_1)^{-k_1}(eN/k_2)^{-k_2}2^{-M}\gamma_0/(4N^{2})$. For the given
$(k_1,k_2)$ there are $\binom{N}{k_1}\binom{N}{k_2}2^{M}$ tuples, since there are
exactly $2^{M}$ functions $[M]\to\bits$, and $\binom Nk\le(eN/k)^{k}$
[CFHS, Lemma~4.3]; all three exponential factors cancel, leaving
$\gamma_0/(2N^{2})$ per pair and $\gamma_0/2\le\gamma_0$ over the $N^{2}$ pairs.
For the reformulation, every subset of $[M]$ is $\theta^{-1}(1)$ for exactly one
$\theta$, and the displayed quantity is
$\bigl|\nu_R(\theta^{-1}(1))-\Unif([M])(\theta^{-1}(1))\bigr|$.
\end{proof}

Neither a conditioning nor a revealed set appears. The bound holds for every
$\theta$ \emph{fixed in advance}, and the test the observer induces --- which does
depend on $f$ --- may then be substituted, because a statement ``for all $\theta$,
$\Phi(f,\theta)\le t$'' may be instantiated at any $\theta$ whatsoever. That single
move dispenses with Definition~\ref{def:rr}, Lemma~\ref{lem:rr} and
Lemma~\ref{lem:mass} altogether, at the price of $M\ln2$ inside $C_1$ where
Lemma~\ref{lem:disc} has $(\sigma+2)\ln2$: the leakage leaves the budget and the
alphabet takes its place.

\begin{remark}[the two routes, and where one source would break both]\label{rem:one}
Lemmas~\ref{lem:disc} and~\ref{lem:disc2} are the same Hoeffding estimate with two
answers to one question: how is a test that depends on $f$ prevented from
conspiring with the very values it is tested against? The first pins the test down
by revealing the cells that determine it, paying $\Pi(\mathcal S)$ for those cells
and only $|\mathcal Z|<2^{\sigma+2}$ for the union; the second declines to pin it
down at all, paying $2^{M}$ for the union and nothing for revealed cells. Neither
dominates, and \S\ref{sec:extraction} carries both. Lemma~\ref{lem:disc2} is not
lossy for what it bounds: by its reformulation the supremum over $\theta$ is a
statistical distance, whose expectation for random $f$ is of order
$\sqrt{M/(k_1k_2)}$ when $M\le k_1k_2$, while
$t_{C_1}\ge\sqrt{M\ln2/(2k_1k_2)}$; Proposition~\ref{prop:F} makes this a proof.
So no smaller test class, chaining argument or sharper tail bound can improve it.

The statement document demands that any proof fail visibly when the two sources
are replaced by one. Both routes do, and not in the size of $t$: at $k_1=1,k_2=N$
one gets $t_{C_0}(1,N)^{2}=(L+N+C_0)/(2N)\to\tfrac12$, so $t_{C_0}\to1/\sqrt2$,
neither large nor vacuous --- and in any case $k_1=1$ models a \emph{deterministic
first coordinate}, not a single source. The failure is earlier and total. A single
source emits a point of $[N]^{2}$ whose posterior given the oracle and the leakage
need not factor at all; Lemma~\ref{lem:prod}(1) is then unavailable, a cell carries
mass up to $\delta$ rather than $m_1m_2$, and there is no product measure for
Lemma~\ref{lem:flat} to flatten onto a rectangle. The union would then range over
$\binom{N^{2}}{k}$ arbitrary $k$-subsets rather than over rectangles, forcing
$t^{2}\ge\frac12\ln(eN^{2}/k)\ge\frac12$ and a vacuous conclusion; and
Lemma~\ref{lem:prod}(4), whose Cauchy--Schwarz step is applied across the two
coordinates, has nothing to be applied to. Lemmas~\ref{lem:flat},
\ref{lem:disc} and~\ref{lem:disc2} have no input and the argument does not begin.
\end{remark}

\section{The observer as a revealing rule}

\begin{definition}[challenge resolution]\label{def:res}
$D$ has \emph{challenge resolution $M'$} if there is $h:[M]\to[M']$ such that for
\textbf{every coin string $\rho$ of $D$}, every $f$, every $\zeta$, and all
$v,v'$ with $h(v)=h(v')$, the runs $D^{f}_{\rho}(v,\zeta)$ and
$D^{f}_{\rho}(v',\zeta)$ issue the same sequence of query \emph{positions}; the
output bit is unconstrained. Every $D$ has resolution $M$; resolution $1$ means
that, for each fixed coin string, the query positions ignore the challenge value.
We may take $h$ surjective without loss: replacing $[M']$ by $\mathrm{image}(h)$
issues fewer runs below and only shrinks the budget.
\end{definition}

The coin quantifier is part of the definition, not a technicality. The weaker
condition that the query positions be challenge-independent \emph{in law} defines
a strictly larger class: at $N=M=2$, $q=1$, the observer that flips $\rho$ and
queries $(1,1)$ when $v\oplus\rho=1$ and $(2,2)$ otherwise has query position
uniform on $\{(1,1),(2,2)\}$ for both challenge values, yet each of its two coin
fixings has resolution $2$. Such an observer steers its probes by the challenge
once its coins are fixed --- exactly the behaviour Barrier~1 leaves open --- and
is not in the class $\kna$ ranges over.

\begin{lemma}[the observer as a revealing rule]\label{lem:rr}
Let $D$ be deterministic, $q$-query, challenge-oblivious, of resolution $M'$.
Then the procedures ``$\mathcal A_\zeta$: for $j=1,\dots,M'$ in order, run
$D^{f}(v_j,\zeta)$ for a fixed $v_j\in h^{-1}(j)$ --- which exists, $h$ being
surjective --- inspecting each queried cell not already known'' form a revealing
rule of budget $s=\min(qM',N^{2})$, with $\theta_{\zeta,f}(v):=D^{f}(v,\zeta)$.
\end{lemma}

\begin{proof}
$\mathcal A_\zeta$ inspects one coordinate at a time, the next being a function
of $\zeta$ and the values already seen: the loop and the representatives are
fixed in advance, and inside run $j$ the next position is a function of
$(v_j,\zeta)$ and previous answers. Each of the $M'$ runs issues at most $q$
positions, so at most $qM'$ cells, and at most $N^{2}$ exist. Finally, for
arbitrary $v$ with $j=h(v)$, the run $D^{f}(v,\zeta)$ issues exactly the
positions run $j$ issued, all inspected and recorded, so its whole execution ---
hence its output bit --- is reconstructible from
$(v,\zeta,\mathcal S,f|_{\mathcal S})$. Therefore the entire function
$\theta_{\zeta,f}$, and $p_{\theta_{\zeta,f}}$, is determined by
$(\zeta,\mathcal S,f|_{\mathcal S})$.
\end{proof}

The resolution parameter is the whole difference between the settled and the open
case. Lemma~\ref{lem:disc} needs the test \emph{as a function on all of} $[M]$,
because the reference $p_\theta$ averages over the alphabet; so every value the
observer might receive must have its query path revealed in advance, and an
observer that steers by $v$ forces $M'=M$.

\begin{lemma}[product mass of the revealed set]\label{lem:mass}
For a revealing rule of budget $s\ge1$,
$\Exp\bigl[\min(\Pi_{f,\zeta}(\mathcal S(f,\zeta)),1)\bigr]\le
\mu(s):=\min\bigl(s\delta,\ 2(s\delta^{2})^{1/3},\ 1\bigr)$; for $s=0$ the left
side is $0$.
\end{lemma}

\begin{proof}
By Lemma~\ref{lem:prod}(1),
$\Pi(\mathcal S)=\sum_{\vu\in\mathcal S}\pi^1(u_1)\pi^2(u_2)\le s\,m_1m_2$.
Taking expectations and using Lemma~\ref{lem:prod}(3) gives the arm $s\delta$;
the arm $1$ is trivial. For the third, fix $\vartheta\in(0,1]$, bound by
$s\vartheta$ on $\{m_1m_2\le\vartheta\}$ and by $1$ off it, and use
Lemma~\ref{lem:prod}(5):
$\Exp[\min(\Pi(\mathcal S),1)]\le g(\vartheta):=s\vartheta+\delta\vartheta^{-1/2}$.
Then $g'(\vartheta)=s-\tfrac12\delta\vartheta^{-3/2}$ vanishes at
$\vartheta^{*}=(\delta/2s)^{2/3}$, which is $\le1$ since $\delta\le1\le2s$, and
$g''>0$; and
$g(\vartheta^{*})=\bigl(2^{-2/3}+2^{1/3}\bigr)(s\delta^{2})^{1/3}\le2(s\delta^{2})^{1/3}$.
\end{proof}

By Remark~\ref{rem:corr} the bound $s\delta^{2}$ is \emph{not} available, and the
exponent $1/3$ is exactly what the $\sqrt\vartheta$ of Lemma~\ref{lem:prod}(5)
produces. Lemma~\ref{lem:mass} is tight at $s=1$ for that example.

\section{Extraction}\label{sec:extraction}

\begin{lemma}\label{lem:ident}
For deterministic $D$ with associated $\theta_{\zeta,f}$,
$\Prob{\Real=1}=\Exp_{(f,\zeta)}\bigl[\Exp_{\vx\sim\Pi_{f,\zeta}}[\theta_{\zeta,f}(f(\vx))]\bigr]$
and $\Prob{\Real_0=1}=\Exp_{(f,\zeta)}[p_{\theta_{\zeta,f}}]$.
\end{lemma}

\begin{proof}
Condition on $(H,\vz)=(f,\zeta)$. In both experiments $D$'s oracle is $f$, so its
output on challenge $v$ is $\theta_{\zeta,f}(v)$. In $\Real$ the challenge is
$f(\vx)$ with $\vx\sim\Pi_{f,\zeta}$ by Lemma~\ref{lem:prod}; in $\Real_0$ it is
uniform and independent, with $\Exp_{v\sample[M]}[\theta(v)]=p_\theta$.
\end{proof}

So $\Real$ and $\Real_0$ differ in the challenge alone, the oracle being the true
$H$ in both: the whole problem is the one-cell question \emph{is $f(\vx)$
distributed like a fresh uniform value, to a reader who knows $\zeta$ and can
probe $f$ $q$ times?}

\begin{proposition}\label{prop:master}
For every $\gamma_0\in(0,1)$ and every deterministic $q$-query
challenge-oblivious $D$ of resolution $M'$,
\begin{multline*}
  \bigl|\Prob{\Real=1}-\Prob{\Real_0=1}\bigr|\le\gamma_0
  +\min\Bigl\{\sqrt{2\delta(L+C_0)},\ \sqrt{2L\delta}+\delta\sqrt{2C_0}\Bigr\}\\
  +\mu\bigl(\min(qM',N^{2})\bigr),
\end{multline*}
with $L=\ln(eN)$ and $C_0$ as in Lemma~\ref{lem:disc}.
\end{proposition}

\begin{proof}
By Lemma~\ref{lem:ident} and the triangle inequality the left side is at most
$\Exp[\Delta]$ with
$\Delta(f,\zeta):=|\Exp_{\Pi}[\theta(f(\vx))-p_\theta]|$. Put
$G_{f,\zeta}(\vu):=(\theta(f(\vu))-p_\theta)\ind[\vu\notin\mathcal S(f,\zeta)]\in[-1,1]$,
so that, since $|\theta-p_\theta|\le1$,
\[
  \Delta(f,\zeta)\ \le\ \bigl|\Exp_{\Pi_{f,\zeta}}[G_{f,\zeta}]\bigr|
  +\Pi_{f,\zeta}(\mathcal S(f,\zeta)),
\]
the two terms being an exhaustive split of the challenge cell. Lemma~\ref{lem:flat}
bounds the first by the largest $|\Exp_{\Unif(B_1)\otimes\Unif(B_2)}[G]|$ over
$|B_i|=k_i:=\lfloor1/m_i\rfloor$, and for such a rectangle that quantity is
\[
  \frac{1}{k_1k_2}\sum_{\vu\in R}G_{f,\zeta}(\vu)
  =\frac{1}{k_1k_2}\sum_{\vu\in R\setminus\mathcal S(f,\zeta)}\bigl(\theta(f(\vu))-p_\theta\bigr),
\]
verbatim the quantity Lemma~\ref{lem:disc} controls; on the event $\mathcal E$ of
that lemma it is at most $t_{C_0}(k_1,k_2)$. As $\Delta\le1$ always, splitting on
$\mathcal E$ gives
$\Exp[\Delta]\le\Prob{\mathcal E^{c}}+\Exp[t_{C_0}]+\Exp[\min(\Pi(\mathcal S),1)]$,
all terms non-negative so the indicators may be dropped; these are at most
$\gamma_0$, Lemma~\ref{lem:mom} at $c=C_0$, and $\mu(s)$ by Lemmas~\ref{lem:rr}
and~\ref{lem:mass}.
\end{proof}

\begin{theorem}[extraction, challenge-blind query positions]\label{thm:A}
Put
\[
  \Lambda\ :=\ \min\bigl\{\,5\sqrt{\sigma'\delta},\ \ \sqrt{2\delta\ln(eN)}
  +4\delta\sqrt{\sigma'}\,\bigr\}.
\]
Then $\kna(q)\le\Lambda+\mu(q)$ for every $q$, and $\kappa(0)\le\Lambda$ with no
restriction on the observer. The second arm of $\Lambda$ is the smaller once $\delta\le0.85$, and is smaller by
a factor $\Theta(\sqrt\delta)$ on the $\sqrt{\sigma'}$ term as $\delta\to0$.
\end{theorem}

\begin{proof}
Both $\kna(q)$ and $\kappa(0)$ are suprema over observers the statement permits
to randomise, whereas Proposition~\ref{prop:master} is stated for deterministic
$D$; by Lemma~\ref{lem:derand} the suprema are attained on deterministic
observers and $D'$ inherits $D$'s challenge resolution, so restricting to
deterministic resolution-$1$ observers loses nothing. If $\delta=1$ then
$\sqrt{\sigma'\delta}\ge\sqrt2>1$ and there is nothing to prove; else take
$\gamma_0:=\delta$ in Proposition~\ref{prop:master}, so $\log\gamma_0^{-1}\le\log N$
as $\delta\ge1/N$. Resolution $1$ gives $s\le q$, so the last term is $\mu(q)$;
for $q=0$, $s=0$ and every observer has resolution $1$ vacuously. It remains to
bound $\gamma_0$ plus each arm of Lemma~\ref{lem:mom} by the corresponding arm of
$\Lambda$.

\emph{First arm.} $\log\gamma_0^{-1}\le\log N\le\sigma'/2$ gives
$\sqrt{2\delta(L+C_0)}\le\sqrt{12\sigma'\delta}\le3.47\sqrt{\sigma'\delta}$, as in
the numerical estimate of the display, and
$\gamma_0=\delta\le\sqrt{\sigma'\delta}$ because $\delta\le1\le\sigma'$; total at
most $5\sqrt{\sigma'\delta}$.

\emph{Second arm.} $2C_0=2(\sigma+2)\ln2+2\ln4+4\ln N+2\ln\delta^{-1}
\le1.3863\sigma+5.5452+6\ln N$, and $\sigma\le\sigma'$, $5.5452\le2.7726\sigma'$,
$6\ln N\le2.0794\sigma'$ --- the last two from $\sigma'\ge2\log N$ and
$\sigma'\ge2$ --- give $2C_0\le6.238\sigma'$, so
$\delta\sqrt{2C_0}\le2.498\,\delta\sqrt{\sigma'}$; and
$\gamma_0=\delta\le0.708\,\delta\sqrt{\sigma'}$ since $\sqrt{\sigma'}\ge\sqrt2$.
The two together are at most $3.21\,\delta\sqrt{\sigma'}\le4\delta\sqrt{\sigma'}$.
For the comparison, $2\delta\ln(eN)\le1.6932\,\sigma'\delta$ gives
$\sqrt{2\delta\ln(eN)}\le1.302\sqrt{\sigma'\delta}$, so the second arm is at most
$(1.302+4\sqrt\delta)\sqrt{\sigma'\delta}$, below $5\sqrt{\sigma'\delta}$ exactly
when $4\sqrt\delta\le3.698$.
\end{proof}

\begin{corollary}\label{cor:A}
If $q^{+}\delta\le\sigma'$ then $\kna(q)\le6\sqrt{\sigma'q^{+}\delta}$. Outside
that regime $q^{+}\delta>\sigma'\ge2$, so $\sqrt{\sigma'q^{+}\delta}>\sigma'>1$
and the target bound is vacuous.
\end{corollary}

\begin{proof}
$\Lambda\le5\sqrt{\sigma'\delta}\le5\sqrt{\sigma'q^{+}\delta}$; and
$q^{+}\delta\le\sigma'$ gives $(q^{+}\delta)^{2}\le\sigma'q^{+}\delta$, i.e.
$\mu(q)\le q\delta\le q^{+}\delta\le\sqrt{\sigma'q^{+}\delta}$.
\end{proof}

This settles the statement's line that queries whose positions ignore the
challenge value are almost free at every budget: their cost is the additive
$q\delta$, with no multiplier on the leading term and no $\log M$.

\begin{theorem}[extraction, unrestricted observers]\label{thm:B}
$\kappa(q)\le\Lambda+\mu(\min(qM,N^{2}))$ for every $q$, with $\Lambda$ as in
Theorem~\ref{thm:A}.
\end{theorem}

\begin{proof}
By Lemma~\ref{lem:derand} the supremum is attained on deterministic observers;
apply Proposition~\ref{prop:master} to such a $D$ with $M'=M$ and
$\gamma_0=\delta$, as in Theorem~\ref{thm:A}.
\end{proof}

\begin{corollary}\label{cor:B}
$\kappa(q)\le6\sqrt{\sigma'q^{+}\delta}$ whenever $q^{+}\delta\le\sigma'$ and
either $M\le\sqrt{\sigma'/(q\delta)}$ or
$M\le\sigma'^{3/2}\sqrt{q/\delta}/8$. If $M\le\sigma'/\sqrt{8\delta}$ then one
of the two holds for every $q\ge1$.
\end{corollary}

\begin{proof}
$qM\delta\le\sqrt{\sigma'q^{+}\delta}$ follows from $qM^{2}\delta\le\sigma'$, and
$2(qM\delta^{2})^{1/3}\le\sqrt{\sigma'q^{+}\delta}$ from
$8qM\delta^{1/2}\le\sigma'^{3/2}q^{3/2}$, both using $q\le q^{+}$; the leading
$\Lambda\le5\sqrt{\sigma'\delta}\le5\sqrt{\sigma'q^{+}\delta}$. The first arm fails
only if $q>\sigma'/(M^{2}\delta)$ and the second only if
$q<64M^{2}\delta/\sigma'^{3}$; both fail at some $q$ only if
$\sigma'^{4}<64M^{4}\delta^{2}$, i.e.\ $M>\sigma'/\sqrt{8\delta}$.
\end{proof}

\subsection*{A bound that does not mention the query budget}

\begin{theorem}[extraction, unrestricted observers, no query budget]\label{thm:D}
For every $\gamma_0\in(0,1)$ and every challenge-oblivious observer $D$ --- of any
query count, bounded or unbounded, and any challenge resolution ---
\[
  \bigl|\Prob{\Real=1}-\Prob{\Real_0=1}\bigr|\ \le\ \gamma_0
  +\sqrt{2\delta\ln(eN)}+\delta\sqrt{2M\ln2+2\ln\tfrac{4N^{2}}{\gamma_0}}\,.
\]
Consequently $\kappa(q)\le5\sqrt{\sigma'\delta}+2\delta\sqrt M$ for every $q$.
\end{theorem}

\begin{proof}
By Lemma~\ref{lem:derand} we may take $D$ deterministic: that lemma bounds no query
count, and nothing here refers to one. By Lemma~\ref{lem:ident} the left side is at
most $\Exp[\Delta]$ with
$\Delta(f,\zeta)=\bigl|\Exp_{\Pi}[\theta_{\zeta,f}(f(\vx))-p_{\theta_{\zeta,f}}]\bigr|$.
Fix $(f,\zeta)$ and put
$G_{f,\zeta}(\vu):=\theta_{\zeta,f}(f(\vu))-p_{\theta_{\zeta,f}}\in[-1,1]$; by
Lemma~\ref{lem:prod} the law $\Pi_{f,\zeta}$ is a product, so Lemma~\ref{lem:flat}
bounds $\Delta$ by the largest
$\bigl|\frac{1}{k_1k_2}\sum_{\vu\in R}G_{f,\zeta}(\vu)\bigr|$ over $|B_i|=k_i$. On
the event of Lemma~\ref{lem:disc2} --- a property of $f$ alone, of probability at
least $1-\gamma_0$ --- that maximum is at most $t_{C_1}(k_1,k_2)$: the pair
$(k_1,k_2)$ lies in $[N]\times[N]$ by Lemma~\ref{lem:flat}, and $\theta_{\zeta,f}$,
although built from $f$, is one particular element of the class that lemma
quantifies over. Since $\Delta\le1$ always,
$\Exp[\Delta]\le\gamma_0+\Exp[t_{C_1}]\le\gamma_0+\sqrt{2L\delta}+\delta\sqrt{2C_1}$
by Lemma~\ref{lem:mom} at $c=C_1$. Nothing above uses the observer's query count,
its resolution, or the revealing apparatus of \S3.

For the consequence, take $\gamma_0:=\delta$; if $\delta=1$ then
$5\sqrt{\sigma'\delta}>1$ and there is nothing to prove. As $\delta\ge1/N$,
\begin{multline*}
  2C_1=2M\ln2+2\ln4+4\ln N+2\ln\delta^{-1}\\
  \le1.3863M+2.7726+6\ln N\le1.3863M+3.4657\,\sigma',
\end{multline*}
using $2.7726\le1.3863\sigma'$ and $6\ln N\le2.0794\sigma'$, so
$\delta\sqrt{2C_1}\le1.1774\,\delta\sqrt M+1.8617\,\delta\sqrt{\sigma'}$. Now
$\delta\sqrt{\sigma'}=\sqrt\delta\sqrt{\sigma'\delta}\le\sqrt{\sigma'\delta}$,
$\gamma_0=\delta\le\sqrt{\sigma'\delta}$, and
$\sqrt{2\delta\ln(eN)}\le1.302\sqrt{\sigma'\delta}$; adding,
$(1+1.302+1.8617)\sqrt{\sigma'\delta}+1.1774\,\delta\sqrt M$.
\end{proof}

For $M=1$ the set $\Fun$ is a singleton, so $\Real$ and $\Real_0$ are the same
experiment and the advantage is $0$; the display is stated for $M\ge2$ only because
$C_1$ is.

\begin{corollary}\label{cor:D}
Let $q\ge1$ and $M\ge4$. If $2\delta\sqrt M<1$ then
$2\delta\sqrt M\le\mu(\min(qM,N^{2}))$, so Theorem~\ref{thm:D} is the smaller of it
and Theorem~\ref{thm:B} wherever either says anything at all. Moreover
$2\delta\sqrt M\le\sqrt{\sigma'q^{+}\delta}$ precisely when
$M\le\sigma'q^{+}/(4\delta)$.
\end{corollary}

\begin{proof}
Against $\mu$'s arm $1$: immediate from $2\delta\sqrt M<1$. Against
$2(s\delta^{2})^{1/3}$ with $s\le qM$: $2\delta\sqrt M\le2(qM\delta^{2})^{1/3}$ iff
$\delta\sqrt M\le q$, and $\delta\sqrt M<1/2<q$. Against $s\delta\le qM\delta$: iff
$q\sqrt M\ge2$, which holds for $q\ge1$ and $M\ge4$. The cap $s=N^{2}$ never binds,
since $\delta\ge1/N$ gives $N^{2}\delta\ge N\ge2$ and
$2(N^{2}\delta^{2})^{1/3}=2(N\delta)^{2/3}\ge2$, so $\mu(N^{2})=1$. The last claim
is $4\delta^{2}M\le\sigma'q^{+}\delta$.
\end{proof}

\section{Presampling, and the decomposition}

\begin{lemma}[presampling, oracle-indexed and consistent]\label{lem:P}
Let $P\in\mathbb N$, $\gamma\in(0,1)$. \emph{(For $M=1$, $\Fun$ is a singleton,
so $\Real_0$ and $\Dec_0$ are the same experiment and the bound holds with left
side $0$; the construction divides by $\log M$ and is stated for $M\ge2$.)} There
is a family $\mathcal Y^{P,\gamma}=\{Y_{f,\zeta}\}$ indexed by $f\in\Fun$ and
$\zeta\in\bits^{*}\times\bits^{*}$ --- with $Y_{f,\zeta}$ uniform on $\Fun$ at
every $\zeta$ outside the support of $\vz$, so the family is total, no experiment
reaching such a $\zeta$ --- depending only on $(S_1,S_2,P,\gamma)$, in which every
$Y_{f,\zeta}$ is a $P$-mixture consistent with $f$, indeed a single
$P$-bit-fixing source, such that for every $q$ and every $q$-query observer
taking $\vz$ and an independent uniform value of $[M]$,
\[
  \bigl|\Prob{\Real_0=1}-\Prob{\Dec_0=1}\bigr|\ \le\
  \frac{q(\sigma+2+2\log\gamma^{-1})}{P}+2\gamma .
\]
\end{lemma}

\begin{proof}
\emph{Step 1: deficiency under randomised leakage.} For $\zeta$ with
$\Prob{\vz=\zeta}>0$ let $\mu_\zeta$ be the law of $H$ given $\vz=\zeta$ and
$S_\zeta:=N^{2}\log M-H_\infty(\mu_\zeta)$. Since
$\mu_\zeta(f)=\Prob{H=f}\Prob{\vz=\zeta\mid H=f}/\Prob{\vz=\zeta}\le(|\Fun|\Prob{\vz=\zeta})^{-1}$,
we get $S_\zeta\le\log(1/\Prob{\vz=\zeta})$. With
$\bar S:=\sigma+2+\log\gamma^{-1}$ and $\mathcal B:=\{\zeta:S_\zeta>\bar S\}$,
every $\zeta\in\mathcal B$ has $\Prob{\vz=\zeta}<\gamma2^{-(\sigma+2)}$, so
$\Prob{\vz\in\mathcal B}<|\mathcal Z|\gamma2^{-(\sigma+2)}<\gamma$. This replaces
the source's own expectation bound on the deficiency, which assumes the advice is
a deterministic function of the oracle and fails here.

\emph{Step 2: the family.} Fix $\zeta$ and set
$\delta_\zeta:=(S_\zeta+\log\gamma^{-1})/(P\log M)$. If $\delta_\zeta>1$ then
$(S_\zeta+\log\gamma^{-1})/P>\log M\ge1$, so the asserted bound exceeds
$q\ge1$ and is vacuous for $q\ge1$, while for $q=0$ both experiments give the
observer a uniform challenge and no oracle access; set $Y_{f,\zeta}$ uniform.
(The governing condition $\delta_\zeta>1$ is $q$-free, so the family is still
defined before $q$.) Otherwise apply the decomposition [CDGS, Claim~2] to
$\mu_\zeta$ with parameter $\delta_\zeta$; since
$(S_\zeta+\log\gamma^{-1})/(\delta_\zeta\log M)=P$ it yields
$\mu_\zeta=\sum_j\lambda_jX_j+\lambda_{\mathrm{fin}}Y_{\mathrm{fin}}$ with
$\lambda_{\mathrm{fin}}\le\gamma$, each $X_j$ a $(P,1-\delta_\zeta)$-dense
source, all supports pairwise disjoint, and each $X_j$ fixing its set $I_j$,
$|I_j|\le P$, to values that every $f$ in its support takes. Let $Y_j$ be the
corresponding $P$-bit-fixing source and define $Y_{f,\zeta}:=Y_j$ for the unique
$j$ with $f\in\supp X_j$, and uniform otherwise. Consistency is immediate; the
construction reads only $(f,\zeta)$ and the decomposition, which depends only on
$(S_1,S_2,P,\gamma)$. Write $J=J(f,\zeta)$.

\emph{Step 3: one component.} Fix $\zeta\notin\mathcal B$ with $\delta_\zeta\le1$
and condition on $J=j$. Because the decomposition is of $\mu_\zeta$ itself,
$\lambda_jX_j$ is the restriction of $\mu_\zeta$ to $\supp X_j$, so the
conditional law of $H$ is exactly $X_j$; in $\Dec_0$ the oracle is drawn freshly
from $Y_j$. The challenge is uniform and independent, so conditioning on it
leaves a fixed $q$-query adaptive distinguisher between $X_j$ and $Y_j$;
[CDGS, Claim~3] bounds its advantage by
$q\delta_\zeta\log M=q(S_\zeta+\log\gamma^{-1})/P$, and averaging over the
challenge preserves this.

\emph{Step 4: averaging.} Averaging over $j$ and bounding the residue event, of
conditional probability at most $\gamma$, by $1$; then averaging over $\zeta$,
bounding $\mathcal B$'s contribution by $\Prob{\vz\in\mathcal B}<\gamma$ and
using $S_\zeta\le\bar S$ off $\mathcal B$, gives the claim.
\end{proof}

\begin{theorem}\label{thm:C}
With $\mathcal Y:=\mathcal Y^{P,\gamma/2}$, for every $q$ and every $q$-query
challenge-oblivious $D$, writing
$\varepsilon(D):=\bigl|\Prob{\Real=1}-\Prob{\Real_0=1}\bigr|$ for that same $D$,
\[
  \Adv_{\mathcal Y,D}\ \le\ \varepsilon(D)
  +\frac{2(\sigma'+\log\gamma^{-1})q}{P}+\gamma+P\delta+q\delta .
\]
Since $\varepsilon(D)\le\kappa(q)$, the same holds with $\kappa(q)$ in place of
$\varepsilon(D)$; but the $D$-relative form is what Theorem~\ref{thm:cprime}
uses, and the inequality $\kappa(q)\le\kna(q)$ is false in general and nowhere
needed.
\end{theorem}

\begin{proof}
Interpolate $\mathsf G_0=\Real$, $\mathsf G_1=\Real_0$, $\mathsf G_2=\Dec_0$,
$\mathsf G_3=\Dec$.

$|\mathsf G_0-\mathsf G_1|=\varepsilon(D)$ by definition, for the fixed $D$ the
theorem is instantiated at; no supremum is taken here.

$|\mathsf G_1-\mathsf G_2|$: Lemma~\ref{lem:P} with slack $\gamma/2$ gives
$q(\sigma+2+2\log(2/\gamma))/P+\gamma$; as $N\ge2$ gives $\sigma+2\le\sigma'$ and
$\sigma'\ge2$, the numerator is at most $2\sigma'+2\log\gamma^{-1}$. For $q\ge P$
that term already exceeds $2\sigma'>1$, so no hypothesis $q<P$ is needed.

$|\mathsf G_2-\mathsf G_3|\le P\delta+q\delta$: draw $H$, $(\vx,\vz)$ and $J$
with fixed set $I_J$; let $H^{\circ}\sim Y_J$ and $U\sample[M]$ be independent of
everything else; let $H^{*}$ be $H^{\circ}$ with its value at $\vx$ overwritten by
$U$ when $\vx\notin I_J$. As $J$ is a function of $(H,\vz)$, $H^{\circ}$ is
independent of $\vx$ given $(H,\vz,J)$; conditioning on $\vx=\vu\notin I_J$, the
source $Y_J$ is uniform and independent at $\vu$, so both $H^{\circ}$ and
$H^{*}$ are distributed as $Y_J$. Running $\mathsf G_3$ with oracle $H^{*}$ and
$\mathsf G_2$ with oracle $H^{\circ}$ and challenge $U$, the two executions
receive the same input on $\{\vx\notin I_J\}$, since there $H^{*}(\vx)=U$, and
their oracles agree off $\vx$. So they diverge only if $\vx\in I_J$ or $D$ queries
$\vx$; being identical up to the step before a first query at $\vx$, that event
has equal probability in both and may be measured in $\mathsf G_2$. The statement
document's lemma \emph{the fixed set misses the challenge} gives
$\Prob{\vx\in I_J}\le P\delta$ --- its proof requires only that $\mathcal Y$
depend on $(S_1,S_2,P,\gamma)$ alone, be indexed by $(f,\zeta)$, and have its
component drawn independently of $\vx$ given $(H,\vz)$, all of which hold, $J$
being a function of $(H,\vz)$. Its lemma \emph{the observer misses the challenge}
gives $\Prob{\vx\in Q}\le q\delta$ in $\mathsf G_2$, its predictor being
simulable there because the challenge is uniform and independent.
\end{proof}

\begin{theorem}[the conjecture, on the region reached]\label{thm:cprime}
Suppose $q^{+}\delta\le\sigma'$ and $P\le\sqrt{\sigma'/\delta}$, and suppose at
least one of
\begin{enumerate}
\item $M\le\sigma'q^{+}/(4\delta)$, with \emph{no} restriction on $D$;
\item $D$ has challenge resolution $1$;
\item the hypothesis of Corollary~\ref{cor:B}, in particular
  $M\le\sigma'/\sqrt{8\delta}$.
\end{enumerate}
Then with $\mathcal Y=\mathcal Y^{P,\gamma/2}$,
\[
  \Adv_{\mathcal Y,D}\ \le\
  \frac{2(\sigma'+\log\gamma^{-1})q^{+}}{P}+8\sqrt{\sigma'q^{+}\delta}+\gamma,
\]
that is, the conjecture holds on that region with $c=2$, $C=8$.
\end{theorem}

\begin{proof}
$\varepsilon(D)\le6\sqrt{\sigma'q^{+}\delta}$ in each branch. Under~1,
Theorem~\ref{thm:D} gives
$\varepsilon(D)\le5\sqrt{\sigma'\delta}+2\delta\sqrt M
\le5\sqrt{\sigma'q^{+}\delta}+\sqrt{\sigma'q^{+}\delta}$, the second term by
Corollary~\ref{cor:D}. Under~2, $D$ has challenge resolution $1$, so
$\varepsilon(D)\le\kna(q)$ and Corollary~\ref{cor:A} applies --- equivalently,
Proposition~\ref{prop:master} bounds $\varepsilon(D)$ directly at $M'=1$. Under~3,
Corollary~\ref{cor:B} bounds $\kappa(q)\ge\varepsilon(D)$. Then
$P\delta\le\sqrt{\sigma'\delta}\le\sqrt{\sigma'q^{+}\delta}$ by the hypothesis on
$P$, and $q\delta\le q^{+}\delta\le\sqrt{\sigma'q^{+}\delta}$ by
$q^{+}\delta\le\sigma'$. Summing, $6+1+1=8$.
\end{proof}

The family is built from $(P,\gamma)$ before $q$ and the bound holds for every
$q$, so the statement's quantifier order is respected; and every $Y_{f,\zeta}$ is
consistent with $f$, so the conclusion is the conjecture and not the weakening
its remark on consistency describes.

Branch~1 is the one that carries no hypothesis on the observer, and it is what
removes from its region the obstruction \S\ref{sec:barriers} calls Barrier~1. It
contains branch~3's clean form whenever $\delta\le1/2$, since then
$\sigma'q^{+}/(4\delta)\ge\sigma'/(4\delta)\ge\sigma'/\sqrt{8\delta}$, and enlarges
it by the factor $\sqrt8\,q^{+}/(4\sqrt\delta)$, which is large exactly when
$\delta$ is small --- the regime the conjecture is about. Branch~2 is \emph{not}
subsumed: it covers every $M$, for that restricted class of observers alone. What
is known is the union of the three.

\section{Combination with the compression bound}

For $q\ge1$, [CFHS, Theorem~3] with $\ell=2$, $2^{-k}=\delta$, $q_D=q$ gives
$\kappa(q)=O\bigl(\log M\cdot(q\delta^{2}(\sigma+2N))^{1/3}\bigr)$. Hence
\[
  \kappa(q)\le5\sqrt{\sigma'\delta}+\min\Bigl\{\mu(\min(qM,N^{2})),\
  O\bigl(\log M\,(q\delta^{2}(\sigma+2N))^{1/3}\bigr)\Bigr\},
\]
and Theorem~\ref{thm:C} converts either arm. The arms are complementary:
Theorem~\ref{thm:B} is smaller when $M\lesssim N\log^{3}M$, the compression arm
when $M$ is very large relative to $N$; the latter also relaxes as $q$ grows,
giving the target shape with no condition on $M$ once
$q\gtrsim\log^{6}M(\sigma+2N)^{2}\delta/\sigma'^{3}$. What neither arm reaches is
$M$ large together with $q$ small. Recorded against that source: its right-hand
side vanishes at $q_D=0$ whereas $\kappa(0)>0$ in general, so it is invoked only
for $q\ge1$, and its constant is not extracted, so this arm is asymptotic while
everything else here is explicit.

\section{The alphabet term is tight}\label{sec:tight}

\begin{proposition}\label{prop:F}
Let $N\ge2$ and $2\le M\le N^{2}/2$. Let $S_1$ and $S_2$ each draw
$x_i\sample[N]$ on their own private coins and output $(x_i,\varepsilon)$. Then the
pair is split, $\sigma_1=\sigma_2=0$, and it is $\delta$-unpredictable with
$\delta=1/N$; and there is a deterministic challenge-oblivious observer $D^{*}$,
making $N^{2}$ queries and of challenge resolution~$1$, with
\[
  \bigl|\Prob{\Real=1}-\Prob{\Real_0=1}\bigr|\ \ge\ \frac{\delta\sqrt M}{4\sqrt2}
  \ =\ \frac{\sqrt M}{4\sqrt2\,N}\,.
\]
\end{proposition}

\begin{proof}
\emph{Parameters.} The coins are independent, so the pair is split. Here
$\vz=(\varepsilon,\varepsilon)$ always, and each $x_i$ is uniform on $[N]$ and
independent of $H$ and of $x_{3-i}$. A predictor is an arbitrary, possibly
randomised function of $(\vz,H)$, so its output is independent of $x_i$ and
$\Prob{\mathsf P^{H}(\vz)=x_i}=1/N$ for every $\mathsf P$: the pair is
$\delta$-unpredictable exactly for $\delta=1/N$, the least value the statement
permits.

\emph{The observer.} $D^{*}$ queries all $N^{2}$ cells, so holds $f$; it forms
$\nu(v):=|f^{-1}(v)|/N^{2}$ and on challenge $v$ outputs $\ind[\nu(v)>1/M]$. It is
challenge-oblivious, deterministic, makes $N^{2}$ queries, and issues the same
query positions for every challenge value, so its resolution is $1$.

\emph{The advantage is an expected statistical distance.} Condition on $H=f$. In
$\Real$ the point $\vx$ is uniform on $[N]^{2}$ and independent of $f$, so the
challenge $f(\vx)$ has law exactly $\nu$; in $\Real_0$ it is uniform on $[M]$.
Hence
\[
  \Prob{\Real=1\mid f}-\Prob{\Real_0=1\mid f}
  =\sum_{v\in[M]}\Bigl(\nu(v)-\tfrac1M\Bigr)^{+}
  =\mathrm{SD}\bigl(\nu,\Unif([M])\bigr)\ \ge\ 0,
\]
so no cancellation occurs on averaging and the advantage is
$\Exp_f[\mathrm{SD}(\nu,\Unif([M]))]$.

\emph{A fourth-moment lower bound.} Let $Y=\sum_{i=1}^{n}Y_i$ with $Y_i$ i.i.d.,
$\Exp Y_1=0$, $|Y_1|\le1$, and $s^{2}:=\Exp Y^{2}=n\Exp Y_1^{2}$; we claim
$\Exp|Y|\ge s/2$ whenever $s\ge1$. Discarding the terms carrying an index to an odd
power, which vanish by independence and centring,
$\Exp Y^{4}=n\Exp Y_1^{4}+3n(n-1)(\Exp Y_1^{2})^{2}\le n\Exp Y_1^{2}+3n^{2}(\Exp Y_1^{2})^{2}
=s^{2}+3s^{4}$, using $\Exp Y_1^{4}\le\Exp Y_1^{2}$, valid as $|Y_1|\le1$. H\"older
with exponents $3/2$ and $3$ gives
$s^{2}=\Exp[|Y|^{2/3}|Y|^{4/3}]\le(\Exp|Y|)^{2/3}(\Exp Y^{4})^{1/3}$, so
$\Exp|Y|\ge s^{3}/\sqrt{s^{2}+3s^{4}}\ge s^{3}/(2s^{2})=s/2$ for $s\ge1$.

\emph{Assembling.} Put $n:=N^{2}$, $p:=1/M$ and $X_v:=|f^{-1}(v)|$. Over uniform
$f$ the indicators $\ind[f(\vu)=v]$ are i.i.d.\ Bernoulli$(p)$, so
$X_v\sim\mathrm{Bin}(n,p)$ and each centred indicator lies in $[-p,1-p]$, of
absolute value at most $1$. Since
$\mathrm{SD}(\nu,\Unif([M]))=\frac{1}{2n}\sum_v|X_v-np|$ and the $X_v$ are
identically distributed, $\Exp_f[\mathrm{SD}]=\frac{M}{2n}\Exp|X-np|$. Here
$s^{2}=np(1-p)=\frac{N^{2}}{M}(1-\frac1M)\ge\frac{N^{2}}{2M}$ as $M\ge2$, and
$s^{2}\ge1$ as $M\le N^{2}/2$; so $\Exp|X-np|\ge\frac{N}{2\sqrt{2M}}$ and
$\Exp_f[\mathrm{SD}]\ge\frac{M}{2N^{2}}\cdot\frac{N}{2\sqrt{2M}}=\frac{\sqrt M}{4\sqrt2N}$.
\end{proof}

Three things follow. First, Theorem~\ref{thm:D} is tight in $M$ to within a factor
$8\sqrt2<12$, so \textbf{no bound free of $q$ can improve its alphabet term} and
Remark~\ref{rem:one}'s claim that Lemma~\ref{lem:disc2} is not lossy is proved.
Second, $D^{*}$ has resolution $1$, so the lower bound applies to $\kna$ as much as
to $\kappa$: whatever remains at large $M$ is \emph{not} the phenomenon Barrier~1
is named for. Consistently, Theorem~\ref{thm:A} says nothing here, $\mu(N^{2})$
being $1$. Third, at $M=\lfloor N^{2}/2\rfloor$ the advantage is at least
$\tfrac18$ while $\sqrt{\sigma'\delta}=\sqrt{2\log N/N}\to0$, so no bound
$\kappa(q)\le h(\sigma',\delta,N)$ free of both $M$ and $q$ can hold: the
statement's $q^{+}$ is necessary. Nothing is refuted --- at those parameters the
target $C\sqrt{\sigma'q^{+}\delta}\ge C\sqrt{2N\log N}$ is itself vacuous.

\section{What is not proved}\label{sec:barriers}

\paragraph{Barrier 1: challenge-steered probes, at $M$ beyond
$\sigma'q^{+}/(4\delta)$.} Uncovered: observers of resolution above $1$ at
$M>\sigma'q^{+}/(4\delta)$, with $q$ in the window left by Corollary~\ref{cor:B}
and the compression arm --- that is, $M$ large together with $q$ \emph{small}. The
mechanism on the revealing route is unchanged and is not slack in the counting.
Lemma~\ref{lem:disc} compares $\Exp_\Pi[\theta(f(\vx))]$ with $p_\theta$, an
average over \emph{all} of $[M]$, so $\theta$ must be known everywhere; by
Lemma~\ref{lem:rr} an observer that steers its probes by the challenge makes
$\theta(v)$ depend on $q$ cells varying with $v$, so all $qM$ must be revealed
before the test is fixed; and Lemma~\ref{lem:mass} is then tight for what it is
given, because coordinatewise $\delta$-unpredictability does not yield
$\Exp[m_1m_2]\le\delta^{2}$ (Remark~\ref{rem:corr}). An observer whose query
positions are challenge-independent only \emph{in law}, steering once its coins are
fixed, has per-coin resolution above $1$ and sits here rather than in
Theorem~\ref{thm:cprime}'s branch~2.

What Lemma~\ref{lem:disc2} does is decline that route entirely, and it is enough
for every $M$ up to $\sigma'q^{+}/(4\delta)$. Beyond that it stops, and
Proposition~\ref{prop:F} shows it stops for a reason: its alphabet term is within a
factor $12$ of the truth, so the residue cannot be reached by any $q$-free
estimate. The sharp question is therefore whether
$\kappa(q)\le O(\sqrt{\sigma'\delta})+\mu(q)$ holds for unrestricted observers ---
whether the revealing route's arm can be made $M$-free. Both arms in hand are
individually tight, $\mu$ at $s=1$ on Remark~\ref{rem:corr}'s family and
$2\delta\sqrt M$ on Proposition~\ref{prop:F}'s, and nothing here claims their
minimum is. Two exits on the revealing route, neither taken: replace the reference
$p_\theta$ by $\theta(f'(\vx))$ for an independent copy $f'$ of $f$ off a single
cell, cutting the budget from $qM$ to $q$ at the price of proving
Lemma~\ref{lem:disc} for a random, $f$-dependent reference; or bound
$\Pi(\mathcal S)$ for \emph{reachable} query sets --- unions of $M$ decision-tree
paths --- rather than arbitrary $qM$-sets. An interpolation between the two routes,
revealing the runs for $v$ in a subset $T\subseteq[M]$ and union-bounding over the
$2^{M-|T|}$ completions, costs $\mu(q|T|)+\delta\sqrt{2(C_0+(M-|T|)\ln2)}$ and
yields nothing beyond the endpoints: driving the first term below $1$ forces
$|T|\lesssim1/(q\delta^{2})$, leaving $M-|T|=M(1-o(1))$ whenever
$\delta\sqrt M\gtrsim1$.

\paragraph{Barrier 2: large $P$ with growing $q$.} Theorem~\ref{thm:C} loses
$P\delta$ at $\mathsf G_2\to\mathsf G_3$, capping the range at
$P\le\sqrt{\sigma'/\delta}$; the statement document names this as a separate open
problem and nothing here improves it. The slack is identifiable: on
$\{\vx\in I_J\}$, which is what costs $P\delta$, consistency makes the decomposed
challenge $H^{*}(\vx)$ equal the real challenge $H(\vx)$, so the event does not
harm $\Adv_{\mathcal Y,D}$ at all --- only the route through the uniform-challenge
worlds. A proof comparing $\mathsf G_3$ with $\mathsf G_0$ directly would not pay
it.

\section{External results used}

\begin{itemize}
\item \textbf{[CDGS, Claim~2]} (Coretti, Dodis, Guo, Steinberger, \emph{Random
  Oracles and Non-Uniformity}, ePrint 2017/937, App.~A). \emph{A distribution on
  $[M]^{n}$ with min-entropy deficiency $S$ is $\gamma$-close to a convex
  combination of finitely many $(P',1-\delta)$-dense sources,
  $P'=(S+\log1/\gamma)/(\delta\log M)$.} The paper states this for a uniform
  variable conditioned on a deterministic function; its proof uses only the
  min-entropy deficiency, and the form above is what that proof establishes. The
  proof's structure --- conditioning on $\{Y_I=y_I\}$ and recursing on
  $\{Y_I\ne y_I\}$ --- also supplies the disjointness of the components'
  supports, the fact that fixed values are values the sampled $f$ takes, and,
  via its recursion invariant, that the decomposition is of the conditioned
  distribution itself, so each component is its restriction and the residue's
  mass is at most $\gamma$ measured under it and not under the uniform
  distribution. All three are used in Lemma~\ref{lem:P}.
\item \textbf{[CDGS, Claim~3]} (same paper). \emph{For a $(P',1-\delta)$-dense
  source $X'$ and its corresponding $P'$-bit-fixing $Y'$, and any adaptive
  $T$-query distinguisher, $|\Prob{D^{X'}=1}-\Prob{D^{Y'}=1}|\le T\delta\log M$.}
\item \textbf{[CFHS, Lemma~3]} (Coretti, Farshim, Harasser, Southern,
  \emph{Multi-Source Randomness Extraction and Generation in the ROM}, ePrint
  2025/1258). \emph{For $k\in\mathbb R_{>0}$ with $2^{k}\in\mathbb N$, every
  $k$-source is a convex combination of flat $k$-sources.}
\item \textbf{[CFHS, Lemma~4.3]} \emph{$(n/k)^{k}\le\binom nk\le(en/k)^{k}$.}
\item \textbf{[CFHS, Theorem~3]} \emph{For $\ell$ sources each
  $(N^{\ell},2^{-k})$-unpredictable with unbounded oracle access and a
  $q_D$-query distinguisher,
  $\Adv^{\mathrm{mse}}=O(\ell\log M\sqrt[\ell+1]{q_D2^{-k\ell}(\sigma+\ell N)})$.}
  Used only for the compression arm.
\item \textbf{Hoeffding's inequality}, two-sided, for independent variables with
  ranges of length $c_i$.
\item \textbf{Cauchy--Schwarz, Markov, Jensen and H\"older}, the last at exponents
  $3/2$ and $3$, in Lemmas~\ref{lem:prod}, \ref{lem:mom}, \ref{lem:mass} and
  Proposition~\ref{prop:F}.
\end{itemize}

\section{Review status}\label{sec:review}

This document has \textbf{not} been independently verified, and a reader should
weigh it accordingly.

Two rounds of adversarial review were run, six passes recording a verdict and two
aborting on infrastructure. \textbf{No pass returned clean.} Five findings were
upheld across the two rounds, of which two were serious; both were repaired, and
both were of the same species --- a claim left underdetermined rather than false,
each recoverable from machinery already present. Zero numbered conclusions were
falsified at any point.

Three limitations of that review are material. Every referee, the adjudicator and
the author are the same vendor's models, so the errors that survive are the
correlated ones by construction; both serious findings came from the strongest
model and were missed by weaker ones examining the same lines. No pass satisfied
the campaign's own blindness condition, which requires a process carrying none of
the project context and was unavailable in this environment. And no protocol
document existed, so every referee ran on instructions written by the session
that produced the proof.

The text above is the third revision, extended. Its own tally is empty: nothing in
it has been reviewed, the repair that produced the third revision touched
Definition~\ref{def:res}, on which every resolution argument depends, and the
extension --- Lemma~\ref{lem:prod}(4),(5), Lemma~\ref{lem:mom},
Lemma~\ref{lem:disc2}, Theorem~\ref{thm:D}, Corollary~\ref{cor:D}, branch~1 of
Theorem~\ref{thm:cprime} and Proposition~\ref{prop:F} --- has been through no
verification pass. One in-session smell test, which is explicitly not such a pass
and recorded no verdict in any tally, raised two items, both repaired here and
neither load-bearing: Lemma~\ref{lem:derand} conjoined the two paired experiments
where its proof supplies a maximiser per pair, and the guide's aside on (D)
claimed of a collected bound what holds only of a display. Its load-bearing new
step is a single application of
Cauchy--Schwarz, Lemma~\ref{lem:prod}(4); a reader with time for one check should
spend it there and on the substitution of $\theta_{\zeta,f}$ into
Lemma~\ref{lem:disc2}, where the quantifier order does all the work.

\end{document}
